%\clearpage
\section {$\nueb$ Appearance Methodology}
\label{sec:method}

The significance of $\nueb$ appearance is given as a p-value by introducing the parameter $\beta$ in order to compare the data to two null hypotheses: a) No $\nueb$ appearance ($\beta = 0$) and b) $\nueb$ appearance according to the current best knowledge of the PMNS matrix ($\beta = 1$), where both hypotheses fully account for uncertainties in the values of the oscillation parameters. Two different test statistics are used: rate-only and rate+shape. 

An ensemble of toys are generated for each null hypothesis, randomising nuisance parameters according to the priors in Tab \ref{tab:osc_param_priors_nuebarapp_optimised} and applying statistical fluctuations. In order to constrain the oscillation and systematic model parameters using T2K data from the ``control samples" (neutrino mode single-ring and multi-ring e-like and both neutrino and antineutrino mode single-ring $\mu$-like), the toy distributions are {\color{red} filled} according to the following method, which ensures that the test-statistic distributions are filled with only the most likely toys and maintains correlations between all nuisance parameters in the 4 ``control samples":

\begin{enumerate}
	\item {\color{red}A large number ($>1M$) of toy experiments are produced taking random parameters from the distributions in Tab \ref{tab:osc_param_priors_nuebarapp_optimised}}
	%and with statistical fluctuations applied.
	
	\item The 4 sample (i.e. excluding antineutrino mode single-ring e-like) likelihood of the data given the toy is calculated for each toy using equation \ref{eq:likelihood_4sample}. Then the maximum likelihood, $\lambda_{max}$, is stored.
	
	\item {\color{red} A proposed toy experiment, $T_{proposed}$, is produced taking random parameters from the distributions in Tab \ref{tab:osc_param_priors_nuebarapp_optimised}.}
	%A toy experiment is produced taking random parameters from the distributions in Tab \ref{tab:osc_param_priors_nuebarapp_optimised} to form the proposed distribution, $T_{proposed}$.
	
	\item \label{item:toy_gen_test} The 4 sample likelihood of the data given $T_{proposed}$ is found: {\color{red}$\lambda_{proposed} = \lambda_{control}(Data; Toy)$.} %\color{red}
	
	\item A random number $R$ is drawn uniformly from $\{0:1\}$.
	
	\item If $\frac{\lambda_{proposed}}{\lambda_{max}}$ is greater than $R$, $T_{proposed}$ is accepted and used in the distribution.
	
	\item The procedure is repeated from step 3 until \NToysFakeDataRatePlusShape\ of $T_{proposed}$ toys have been accepted, denoted $T_{accepted}$. The resulting set of $T_{accepted}$ toys are constrained by the likelihood of the ``control sample" data given the toy and only contain the most likely toys. {\color{red} Toys generated by this method will be referred to as data-constrained toys. }
	
	\item The rate-only and rate+shape test-statistics are calculated according to sections \ref{sec:method_rate-only} and \ref{sec:method_rate+shape} for each $T_{accepted}$ to produce test-statistic distributions.
\end{enumerate}

\begin{equation} \label{eq:likelihood_4sample}
\lambda_{control} \left( Data; Toy \right) =
\prod_{{\color{red} s \ne RHC1Re }} \lambda_s \left( Data; Toy \right)
\end{equation}

It is worth noting that this method of filling the toy distributions is statistically equivalent to performing a likelihood ratio fit of the ``control samples" and using the posterior distribution of the nuisance parameters as a prior to extract the systematic parameters for a toy experiment.

At this point, the rate-only and rate+shape analysis methods diverge.



\subsection{Rate-only Analysis Method}
\label{sec:method_rate-only}

The total number of observed events, $N_{obs}$, in the antineutrino mode e-like sample, without using any information from the reconstructed energy and angle distribution, is used to quote the rate-only p-value. 

The expected distribution of antineutrino mode single-ring e-like events, under statistical and systematic fluctuations, is produced using \NToysFakeDataRateOnly\ (throws of nuisance parameters) x \NStatThrows\ (throws of statistical fluctuations) toy data sets to ensure the nuisance parameter space was scanned effectively. In this method, the toy experiments used for the \NToysFakeDataRateOnly\ nuisance parameter throws are the sampled toys described in section \ref{sec:method}. For each of these toys, \NStatThrows\ statistical throws are performed using a Poisson distribution with mean equal to the number of antineutrino mode single-ring e-like events in the toy. The result of each statistical throw, $N_{exp}$, is stored in {\color{red}a rate-only test-statistic distribution for each toy, $f_{itoy} \left( N_{exp} \right)$.} %\color{red}
%The result of each statistical throw, $N_{exp}$, is stored in the rate-only test-statistic distribution, $f \left( N_{exp} \right)$. 


After performing statistical throws, a p-value is calculated for {\color{red} each $f_{itoy} \left( N_{exp} \right)$. } For the $\beta = 0$ null hypothesis, a right-tailed p-value is used (see equation \ref{eq:pval_right}), while for the $\beta = 1$ null hypothesis, a left-tailed p-value is used (see equation \ref{eq:pval_left}). {\color{red}For each null hypothesis, the mean of the rate-only p-values (averaged over each toy) is quoted as the result. This is equivalent to storing all statistical throws for all toys in a single distribution and calculating a single p-value from this, but is done to increase computational efficiency. } %\color{red}
%After performing statistical throws, a p-value is calculated for $f \left( N_{exp} \right)$. For the $\beta = 0$ null hypothesis, a right-tailed p-value is used (see equation \ref{eq:pval_right}), while for the $\beta = 1$ null hypothesis, a left-tailed p-value is used (see equation \ref{eq:pval_left}). 

A right-tailed p-value is used for the $\beta = 0$ hypothesis as the $\beta = 0$ test-statistic distribution is composed solely of background events. The p-value is the probability of obtaining a result as extreme or more extreme than the observed result under the null hypothesis, so using a right-tailed p-value gives the probability of observing a fluctuation under the $\beta = 0$ hypothesis in the background with equal to or more events than observed in the antineutrino mode single-ring e-like sample. An upwards fluctuation such as this would be less consistent with coming from the expected number of events from a sample containing only background events. Therefore, if this p-value is lower than some threshold (usually set to $\SI{5}{\sigma}$), it could be claimed that the observation is incompatible with having come from the null hypothesis and the hypothesis is rejected. 

Similarly, a left-tailed p-value is used for the $\beta = 1$ hypothesis as the $\beta = 1$ test-statistic distribution is composed of both background and signal events, so the left-tailed p-value gives the probability of observing a fluctuation under the $\beta = 1$ hypothesis with equal to or fewer events than were observed in the antineutrino mode single-ring e-like sample. A downwards fluctuation such as this would be less consistent with coming from the expected number of events from a sample containing both background and signal events.

\begin{equation} \label{eq:pval_right}
\textrm{p-value} \textrm{(right tail)} = 
\frac{\int_{N_{obs}}^{+\infty} f \left( N_{exp} \right) dN_{exp}}{ \int_{-\infty}^{+\infty} f \left( N_{exp} \right) dN_{exp}}
%\textrm{p-value}_{itoy} \textrm{(right tail)} = 
%\frac{\int_{N_{obs}}^{+\infty} f_{itoy} \left( N_{exp} \right) dN_{exp}}{ %\int_{-\infty}^{+\infty} f_{itoy} \left( N_{exp} \right) dN_{exp}}
\end{equation}

\begin{equation} \label{eq:pval_left}
\textrm{p-value} \textrm{(left tail)} = 
\frac{\int^{N_{obs}}_{-\infty} f \left( N_{exp} \right) dN_{exp}}{ \int_{-\infty}^{+\infty} f \left( N_{exp} \right) dN_{exp}}
%\textrm{p-value}_{itoy} \textrm{(left tail)} = 
%\frac{\int^{N_{obs}}_{-\infty} f_{itoy} \left( N_{exp} \right) dN_{exp}}{ %\int_{-\infty}^{+\infty} f_{itoy} \left( N_{exp} \right) dN_{exp}}
\end{equation}

%\begin{equation} \label{eq:pval_mean}
%\textrm{p-value} = \frac{1}{n_{toys}} \sum_{itoy = 0}^{n_{toys}-1} \textrm{p-value}_{itoy}
%\end{equation}

As the distribution of the statistic is populated with such a large sample, the statistical uncertainty on the p-value is very small, of the order of $10^{-3}$.





\subsection{Rate+shape Analysis Method}
\label{sec:method_rate+shape}

{\color{red}
In calculating the rate+shape test statistic, the marginal likelihood is used and is calculated according to the following method.


\subsubsection{Calculation of the Marginal Likelihood} \label{sec:marg}

The marginal likelihood is computed by considering an ensemble of MC predictions, generated based on priors on systematic parameters, as in \cite{OurRun18CC1piJointAnalysis3Flavour}, and priors on oscillation parameters from Tab \ref{tab:osc_param_priors_nuebarapp_optimised}. Together, the oscillation and systematic parameters, including the mass hierarchy, are considered to be nuisance parameters. The dependence of the likelihood on the nuisance parameters is eliminated during the marginalisation process so that the likelihood is only dependant on $\beta$. The marginal likelihood is then calculated as the weighted mean of the antineutrino mode single-ring e-like sample likelihood evaluated at $N_{marg}$ throws of all nuisance parameters, ${\bf a}$:

\begin{equation} \label{eq:marginal_likelihood}
\lambda_{marg} \left( \beta \right) = \frac{1}{N_{marg}} \sum_{j=1}^{N_{marg}} 
\lambda_{RHC1Re} (T_{accepted}; {\bf a}_{j},\beta) \times
\lambda_{control}(Data; {\bf a}_{j},\beta)
\end{equation}

Where $\lambda_{RHC1Re}(T_{accepted}; {\bf a}_{j},\beta)$ denotes the likelihood of each $T_{accepted}$ in the antineutrino mode single-ring e-like sample for the $j^{\text{th}}$ set of nuisance parameter throws, ${\bf a}_{j}$. This is weighted by the control sample likelihood of the data given the $j^{\text{th}}$ set of nuisance throws, $\lambda_{control}(Data; {\bf a}_{j},\beta)$. Weighting the likelihood in this way is equivalent to using the control sample data as a prior constraint when throwing the nuisance parameters.
%is equivalent to having thrown the nuisance parameters using a constraint on the nuisance parameters from the control sample data as as a prior.

In the limit of large statistics, $-2\ln \lambda$ asymptotically converges to a $\chi^2$ distribution \cite{WilksTheorem}. From this point onwards, the marginal log-likelihood will be instead written as $\chi^2_{marg}$.



\begin{equation} \label{eq:marginal_chisq}
\chi^2_{marg} \left( \beta \right) 
= -2 \; \textrm{ln} \; \lambda_{marg}
\left( \beta \right)
\end{equation}

In this analysis, \NToysMarg\ nuisance parameter throws are used to determine the marginal likelihood in order to explore the whole space of the systematic parameters. More throws are used compared to the joint fit analysis, because in this analysis the mass hierarchy, $\delta_{CP}$ and all other oscillation parameters are marginalised at the same time whereas they were not in the joint analysis. Furthermore these parameters have a large impact on the $\numu \rightarrow \nue$ component which is a significant background in the null hypotheses.



\subsubsection{Calculation of the Test Statistic and p-value}

} %\color{red}

The rate+shape analysis method uses log-likelihood ratio between the $\beta = 0$ and $\beta = 1$ $\nueb$ appearance null hypotheses as the test-statistic:

\begin{equation} \label{eq:Delta_chisq}
\Delta \chi^2 = \chi_{marg}^2 \left( \beta = 0 \right) - \chi_{marg}^2 \left( \beta = 1 \right)
\end{equation}


This test statistic simply compares the oscillation with the non-oscillation hypothesis but, unlike the $\Delta \chi^2$ values given in the 3-flavour joint analysis \cite{OurRun18CC1piJointAnalysis3Flavour}, this does not use the best-fit $\chi^2$ value, as $\beta$ is not fit for specifically. For this reason the distribution of $\Delta \chi^2$ is expected to take both negative and positive values, where $\Delta \chi^2 < 0\ (> 0)$ corresponds to the spectra being more consistent with the $\beta = 0\ (1)$ hypothesis than $\beta = 1\ (0)$. It is worth noting that the parameter $\beta$ is not fit, but is fixed either to 0 or 1. %It is also worth noting that marginal $\chi^2$ is calculated by using all 5 Super-K samples.

The $\Delta \chi^2$ distribution is produced by first generating \NToysFakeDataRatePlusShape\ toys produced using the method described in section \ref{sec:method}. The toys are then marginalised using \NToysMarg\ marginalisation toys generated from the priors in Tab \ref{tab:osc_param_priors_nuebarapp_optimised} and $\chi^2_{marg} \left( \beta \right)$ is calculated according to equation \ref{eq:marginal_chisq}. This incorporates the uncertainty on nuisance parameters on the measurement and distributions of the test-statistic.
%It is worth noting that this takes into account only systematic fluctuations not statistical fluctuations.

Once a $\Delta \chi^2$ value is calculated for each toy, it is stored in the test-statistic distribution, $f \left( \Delta \chi^2 \right)$, which is integrated to obtain p-values, for example the right-tailed p-value for the $\beta = 0$ null hypothesis shown in equation \ref{eq:pval_rate+shape_right}. %Although this is in contrast to the rate-only analysis, where there is one test-statistic distribution per toy, both methods for calculating the p-value are mathematically identical. The method used to calculate the rate-only p-value is used in order to produce distributions of the toy p-value.

%Once a $\Delta \chi^2$ value is calculated for each toy, it is stored in a single test-statistic distribution

\begin{equation} \label{eq:pval_rate+shape_right}
\textrm{p-value} \textrm{(right tail)} = 
\frac{\int_{\Delta \chi^2_{obs}}^{+\infty} f \left( \Delta \chi^2 \right) d\Delta \chi^2}{\int_{-\infty}^{+\infty} f \left( \Delta \chi^2 \right) d\Delta \chi^2}
\end{equation}


{\color{red}
	\subsection{Differences Between the Rate-only and Rate+shape Analysis Methods}
} %\color{red}

The number of statistical throws from each toy data set is different between the rate-only and rate+shape analyses. The rate+shape analysis {\color{red} uses no single statistical fluctuations when} generating each fake-data toy. %The reason for this is because the limiting factor in the rate+shape analysis is the calculation of the marginal likelihood.

It is worth pointing out that sometimes the rate-only p-value can be smaller than the rate+shape. This can be due to both the discreteness of the rate-only p-value or to the shape information (which the rate-only analysis does not have access to) which can discredit the $\nueb$ appearance hypothesis. For this reason it is important to quote both the p-values as our sensitivity.\\


{\color{red}
	\subsection{Comparison to the P-theta Analysis}
} %\color{red}

Both the rate-only and rate+shape analyses were validated against what was found by the P-theta group. {\color{red} In the Run 1-9c analysis for} the rate-only analysis we see good agreement between the two analyses, as expected. For the rate+shape analysis, the difference between P-theta and VALOR is mainly due to the e-like samples. For VALOR the e-like samples are binned in $E_{reco}-\theta$ analysis templates whereas for P-theta they are binned in $p-\theta$. When the P-theta group runs the analysis using $E_{reco}-\theta$, the results of VALOR and P-theta agree reasonably well {\color{red} (see Ref \cite{nuebarapp_valor_plenary_run1-9c} for the Run 1-9c analysis). For the full Run 1-9 analysis, good agreement between VALOR and P-theta is found, as shown in appendix \ref{app:compare_valor_ptheta}.}

{\color{red}
In the Run 1-9c analysis, the remaining differences were thought to be due to the different toy generation methods used by VALOR and P-theta: Instead of sampling the toy parameter space around the maximum control sample likelihood to accept/reject toys, the P-theta analysis accepts all toys and applies the control sample likelihood as a weight when calculating p-values. This affect is studied in appendix \ref{app:sampling_vs_reweighting} using full Run 1-9 statistics.
} %\color{red}

In addition, the P-theta group performed fake-data studies to assess the impact of uncertainties in the neutrino-nucleus interaction model and found no evidence of any significant bias \cite{PthetaR19cNuebarAppearance}.\\


\subsection{Statistical Uncertainty on the p-values}
The statistical uncertainty on the p-value is calculated assuming that the number of entries in each bin of the expected distribution of the test statistic follow a Poisson distribution. The p-value uncertainty is then derived from the statistical uncertainty of the population. As the p-value is calculated from the fraction of trials which fall to the right of the data statistic, the p-value is expected to be binomially distributed. Given a sufficiently large number of trials ($N$), the p-value is expected to take a gaussian distribution with a mean of the p-value and a standard deviation of $\sqrt{p(1-p)/N}$.

The uncertainty is calculated from the first order Taylor series expansion assuming no correlation between bins in the test-statistic distribution and the number of events in each bin follows a Poisson distribution. The calculation of the uncertainty summing over all toy experiments, for the right sided rate-only p-value, is shown as an example and is as follows:

The p-value calculation shown in equation \ref{eq:pval_right} is written in terms of an integral over the test-statistic distribution. This is re-written as a sum of all bins that contribute to the right-tail of the distribution

\begin{equation} \label{eq:pval_pn_pd}
\textrm{p-value} \textrm{(right tail)} = p =
\frac{\int_{N_{obs}}^{+\infty} f \left( N_{exp} \right) dN_{exp}}{ \int_{-\infty}^{+\infty} f \left( N_{exp} \right) dN_{exp}} =
\frac{\sum_{i=0}^{n(right-tail)} n_{i}}{\sum_{i=0}^{n(total)} n_{i}}
\end{equation}

Where $n_{i}$ is the the number of times each toy, $i$, appears. The total uncertainty as a result of the statistical uncertainty on each sample is given by

\begin{equation} 
\left( \sigma_p \right)^2 = 
\sum_{i} \left( \frac{\partial p}{\partial n_i} \sigma_{n_i} \right)
\end{equation}

Expanding and using the Poisson uncertainty on the toys $\sigma_{n_i} = \sqrt{n_i}$ before taking the limit where each toy appears only once gives

\begin{equation} \label{eq:pval_error}
\left( \sigma_p \right)^2 = 
\frac{p(1-p)}{N_{tot}}
\end{equation}

Where $N_{tot}$ is the integral of the test-statistic distribution. This is the expected result for a binomial process with a success rate $p$.



%The uncertainty in p given uncertainty in events per bin. Considering bins to the right and left of the data statistic independently:

%\begin{equation}
%\left( \sigma_p \right)^2 = 
%\sum_{i}^{n(right-tail)} \left( \frac{p_d - p_n}{p_d^2} \sigma_{n_i} \right)^2 + 
%\sum_{i}^{n(left-tail)} \left( \frac{p_n}{p_d^2} \sigma_{n_i} \right)^2
%\end{equation}

%Where we have used the Poisson uncertainty on the toys $\sigma_{n_i} = \sqrt{n_i}$. Which, when using the definition of $p$, becomes:



\subsection{Bayes Factor} \label{sec:bayes}
As in this analysis the $\Delta \chi^2$ statistic considered is marginalised over all nuisance parameters, the statistic can be interpreted in a Bayesian framework as a Bayes factor. The Bayes factor is another statistic (fully defined by the rate+shape analysis $\Delta \chi^2$ statistic) and gives an indication of which hypothesis gives the best fit to the data. The Bayes factor is defined as the marginal likelihood ratio of the two hypotheses in question.

\begin{equation} \label{eq:bayes_factor}
B_{01} = \frac{e^{-\chi^2_{\beta = 0}/2}}{e^{-\chi^2_{\beta = 1}/2}} = e^{-\frac{1}{2}\Delta \chi^2}
\end{equation}

The Jeffreys scale \cite{JeffreysScale} is used to interpret the Bayes factor to show whether there is a preference for one hypothesis over the other. A more detailed discussion of Bayes factors can be found in \cite{MaCH3NuebarAppearance2015}.


%%% Local Variables:
%%% mode: latex
%%% TeX-master: "../TN327"
%%% End:

